\subsection{FVM 基準手法の定式化と精度解析}
\label{app:fvm_reference_methods}

% ═══════════════════════════════════════════════════════════════════════════════
\subsubsection{\texorpdfstring{PPE 面係数 $a_f = 2/(\rho_L+\rho_R)$ の積分論理と周期境界対応}{PPE 面係数 a\_f = 2/(ρ\_L+ρ\_R) の積分論理と周期境界対応}}
\label{app:fvm_face_coeff}
\label{sec:fvm_face_coeff}
% ═══════════════════════════════════════════════════════════════════════════════
%% ※ 旧 08g_fvm_note.tex の内容を統合

PPE 演算子 $\bnabla\cdot(a\bnabla p)$（$a \equiv 1/\rho$）を FVM で離散化すると，
フェイス面係数として\textbf{調和平均}が自然に現れる：
\begin{equation}
  a_f = \frac{2}{\rho_L + \rho_R}
  \label{eq:af_exact}
\end{equation}
よくある係数取り違えとして，$\rho$ の調和平均の逆数
$(\rho_L+\rho_R)/(2\rho_L\rho_R)$ は式~\eqref{eq:af_exact} ではない．
式~\eqref{eq:af_exact} は，低次 FVM 基準離散化，周期境界の参照行列，
および §~\ref{sec:fvm_ccd_corrector} の面フラックス整合を説明するための
面係数である．本稿の高密度比・低 Weber 数経路では，smoothed-Heaviside
一括 PPE にこの係数だけを入れて閉じるのではなく，
分相 PPE + 圧力ジャンプ条件 + HFE（§~\ref{sec:split_ppe_framework},
§~\ref{sec:field_extension}）へ切り替える．

\subsubsection{FVM 積分論理}

セル $P$（位置 $x_L$）と $E$（位置 $x_R$）の間の等価面係数を FVM 積分で求める．
フラックス $F = a_f(p_E - p_P)/\Delta x$ において，
$a(x)$ が区分的定数（$a_L = 1/\rho_L$，$a_R = 1/\rho_R$）のとき：
\[
  \int_{x_L}^{x_R} \frac{dx}{a(x)}
  = \frac{\Delta x/2}{a_L} + \frac{\Delta x/2}{a_R}
  \quad\Rightarrow\quad
  a_f = \frac{2a_L a_R}{a_L+a_R}
  = \frac{2}{\rho_L+\rho_R}
\]
これは直列抵抗の等価抵抗と同じ論理（調和平均）である．\qed

\subsubsection{FVM PPE 行列の周期境界条件対応}
\label{sec:fvm_periodic}
%% 旧所在: §8.4（08f_ppe_bc.tex）→ 付録 E.3 へ移動

壁面（ノイマン）BC に対する FVM 離散化（式~\eqref{eq:af_exact}）は，
周期境界を持つ計算領域ではそのまま適用できない．
低次 FVM 基準離散化では以下の3点を変更して周期 FVM 行列を組み立てる．

\paragraph{（1）折り返しフェイスの追加}
周期 BC では内部フェイスに加え，ノード $N_x-1$ とノード $0$ を結ぶ
\textbf{折り返しフェイス（wrap face）}を追加する：
\begin{equation}
  a_{N_x - 1/2} = \frac{2}{\rho_{N_x-1} + \rho_0}
  \label{eq:af_wrap}
\end{equation}

\paragraph{（2）境界半格子体積補正の省略}
周期 BC ではすべてのノードが同じ格子幅 $h$ の制御体積を持つため，
壁面 BC の $h/2$ 体積補正は行わない（$a_f / h^2$ を一様に適用）．

\paragraph{（3）ゴーストノードの同一性拘束}
インデックス $N$ ノード（物理的に $0$ の周期像）を
\textbf{同一性拘束行}で置き換える：
\begin{equation}
  p_{\mathrm{ghost}} - p_{\mathrm{src}} = 0
  \quad\Longleftrightarrow\quad
  A[\mathrm{ghost},\,\mathrm{ghost}] = 1,\;
  A[\mathrm{ghost},\,\mathrm{src}]   = -1,\;
  b[\mathrm{ghost}] = 0
  \label{eq:ghost_identity}
\end{equation}
コーナー節点は最初の軸の拘束のみを採用する（重複適用を避けるため）．

この同一性拘束は PPE だけの数値線形代数上の都合ではなく，
周期領域を離散商空間として扱う幾何条件である．
任意の節点場 $q_h$（圧力，速度，$\psi$，曲率補助場）は
式~\eqref{eq:periodic_quotient_space} の $\mathcal{Q}_P$ 上で評価し，
各演算子適用後に周期像面を source 面へ同期する．
とくに面中心 flux から節点速度へ戻す場合，周期軸 $a$ では
壁境界用の片側 endpoint 復元ではなく
\begin{equation}
  u_{i}^{(a)}
  =
  \frac{1}{2}
  \left(
    u_{f,i}^{(a)}
    +
    u_{f,i-1\ \mathrm{mod}\ N_a}^{(a)}
  \right),
  \qquad
  u_{N_a}^{(a)}=u_{0}^{(a)}
  \label{eq:periodic_face_to_node_reconstruction}
\end{equation}
を用いる．
混合境界（例：$S^1\times[0,L_y]$）では，周期軸には
式~\eqref{eq:periodic_face_to_node_reconstruction} を適用し，
wall 軸にのみ片側壁閉包と no-slip/no-penetration 条件を適用する．

\begin{tcolorbox}[warnbox, title={周期 FVM 行列のゲージと比較用途}]
  式~\eqref{eq:ghost_identity} の同一性拘束行と FVM 行が混在した疎行列に対して，
  Krylov 前処理の性能を本稿の検証契約に含めない．低次 FVM 基準値は
  ゲージ固定後の直接法（\texttt{spsolve}）で取得し，反復前処理の
  収束性は主張範囲から外す．
\end{tcolorbox}

% ══════════════════════════════════════════════════════════════════════════════
\subsubsection{CSF（連続表面力）モデルの参照定式化}
\label{app:csf_full}
% ══════════════════════════════════════════════════════════════════════════════

\textbf{連続表面力（CSF: Continuum Surface Force）モデル} \cite{Brackbill1992} は
表面張力（面力）を体積方程式に組み込む標準的手法である．
本付録では，第~\ref{sec:levelset}章§\ref{sec:smoothed_delta} で導入した
平滑化デルタ関数 $\delta_\varepsilon$（式~\eqref{eq:delta_eps_def}，
図~\ref{fig:heaviside_delta}）を用いて CSF モデルを定式化する．

% ───────────────────────────────────────────────────────────────────────────
\subsubsection{界面ディラックデルタと CSF 体積力}
% ────────────────────────────────────────────────────────────────────────────

界面 $\Gamma$ 上に作用する表面張力は，本稿の液相$\to$気相法線規約では
単位面積あたり $-\sigma\kappa_{lg}\hat{\bm{n}}_{lg}$ の法線力である．
これを全空間で定義された体積力 $\bm{f}_\sigma$ として書くには，
界面上だけ値を持つ\textbf{界面ディラックデルタ} $\delta_s$ を用いる：
\begin{equation}
  \bm{f}_\sigma = -\sigma\kappa_{lg}\hat{\bm{n}}_{lg}\,\delta_s(\bm{x})
  \label{eq:csf_dirac_app}
\end{equation}
§\ref{sec:smoothed_delta} で示したように，CLS 変数の勾配
$|\bnabla\psi| \approx \delta_\varepsilon(\phi)$（式~\eqref{eq:grad_psi_delta}）が
$\delta_s$ の滑らかな近似を与え，界面法線は
$\hat{\bm{n}}_{lg} = -\bnabla\psi/|\bnabla\psi|$ で得られる．

% ─────────────────────────────────────────────────────────────────────────────
\subsubsection{CSF 体積力の最終形：\texorpdfstring{$\bm{f}_\sigma = \sigma\kappa_{lg}\bnabla\psi$}{f\_sigma = sigma*kappa_lg*grad(psi)}}
% ────────────────────────────────────────────────────────────────────────────

式~\eqref{eq:grad_psi_delta} の関係を式~\eqref{eq:csf_dirac_app} に代入する：
\begin{equation}
  \bm{f}_\sigma = -\sigma\kappa_{lg}\,\hat{\bm{n}}_{lg}\,\delta_s
  \approx -\sigma\kappa_{lg}\,\frac{-\bnabla\psi}{|\bnabla\psi|}\,|\bnabla\psi|
  \approx \sigma\kappa_{lg}\,\bnabla\psi
  \label{eq:csf_final}
\end{equation}

ただし，この近似は $|\bnabla\phi| \approx 1$（Eikonal 条件）が成立する場合に限り有効である．
CLS 法では再初期化ステップによって $|\bnabla\phi|\approx 1$ を維持し，この近似の誤差を抑制する．
ただし精度は曲率評価，界面厚み，離散演算子の整合性にも依存するため，別途検証で確認する．

$\bnabla\psi$ は気相から液相を向き，界面で最大値をとる．
したがって $\sigma\kappa_{lg}\bnabla\psi$ は界面近傍でのみ作用し，
全空間積分すると Young-Laplace のジャンプ条件 $j_{gl}=p_g-p_l=-\sigma\kappa_{lg}$ を再現する
（以下の検証を参照）．

\medskip
\noindent\textbf{検証：CSF モデルが Young-Laplace ジャンプ条件を再現すること}
\label{sec:csf_young_laplace_app}

CSF 体積力 $\bm{f}_\sigma = \sigma\kappa_{lg}\bnabla\psi$ が界面ジャンプ条件
$j_{gl}\equiv p_{\text{gas}} - p_{\text{liq}} = -\sigma\kappa_{lg}$
を再現することを，界面法線方向の積分によって示す．

平衡状態（速度ゼロ）では運動方程式より
$\bnabla p = \sigma\kappa_{lg}\bnabla\psi$ が成立する．
界面を法線方向 $s$ に沿って積分すると：
\[
  p_{\text{gas}} - p_{\text{liq}}
  = \sigma\kappa_{lg} \int_{\text{liq}}^{\text{gas}} \frac{\partial\psi}{\partial s}\,ds
  = \sigma\kappa_{lg}\,(0 - 1) = -\sigma\kappa_{lg}
\]
\begin{equation}
  p_{\text{gas}} - p_{\text{liq}} = -\sigma\kappa_{lg}
  \label{eq:young-laplace-app}
\end{equation}

% ══════════════════════════════════════════════════════════════════════════════
\subsubsection{Balanced-Force 条件：離散化演算子不一致の誤差解析}
\label{app:balanced_force_taylor}
% ═════════════════════════════════════════════════════════════════════════════

§\ref{sec:bf_operator_mismatch}（式~\eqref{eq:bf_operator_mismatch}）で示した演算子不一致誤差を，離散テイラー展開によって定量的に導出する．

\subsubsection{打ち切り誤差の定量評価}

連続レベルでは $\bnabla p = \kappa\bnabla\psi/We$ の釣り合いが成立している静止液滴を，
異なる差分演算子 $D^{(1)}_{\mathrm{low}}$（低次）と $D^{(1)}_{\mathrm{CCD}}$（CCD，高次）で計算した場合の離散誤差を見る．

$f$ が解析関数のとき，各演算子の打ち切り誤差は：
\begin{align}
  D^{(1)}_{\mathrm{low}} f   &= f' + c_1 h^2 f''' + \Ord{h^4} \\
  D^{(1)}_{\mathrm{CCD}} f   &= f' + c_2 h^6 f^{(7)} + \Ord{h^8}
\end{align}

\textbf{平衡条件 $\nabla p = (\kappa/We)\nabla\psi$ を用いた残差の展開：}
\begin{align}
  D^{(1)}_{\mathrm{CCD}} p - \frac{\kappa}{We} D^{(1)}_{\mathrm{low}}\psi
  &= \Bigl[\nabla p + \Ord{h^6}\Bigr]
     - \frac{\kappa}{We}\Bigl[\nabla\psi + \Ord{h^2}\Bigr] \notag\\
  &= \underbrace{\Bigl(\nabla p - \frac{\kappa}{We}\nabla\psi\Bigr)}_{=\,0\;\text{（平衡条件）}}
     + \Ord{h^6} - \frac{\kappa}{We}\Ord{h^2} \notag\\
  &= \Ord{h^2}
  \label{eq:bf_mismatch_expand}
\end{align}
FD の打ち切り誤差 $O(h^2)$ が残差を支配し（CCD の $O(h^6)$ は無視できる），この $O(h^2)$ の不整合が\textbf{体積力として運動量方程式に残り，寄生流れを駆動し，最終的に計算の爆発を引き起こす}（§\ref{sec:bf_operator_mismatch}）．

\subsubsection{同一演算子（CCD）使用時の誤差低減}

同一の演算子（CCD）を両方に用いると，各項が独立に $O(h^6)$ 精度で連続微分を近似するため（$D_{\mathrm{CCD}}^{(1)} f = \nabla f + O(h^6)$），平衡条件 $\nabla p = (\kappa/We)\nabla\psi$ のもとで：
\[
  D^{(1)}_{\mathrm{CCD}} p - \frac{\kappa}{We} D^{(1)}_{\mathrm{CCD}}\psi
  = \underbrace{\nabla p - \frac{\kappa}{We}\nabla\psi}_{=\,0\;\text{（平衡条件）}} + O(h^6)
  = O(h^6)
\]
となり，寄生流れの\textbf{数値離散化誤差成分}は $O(h^6)$ まで低減される．

\noindent\textbf{（注：演算子の可換性について）}
$D_{\mathrm{CCD}}(\kappa\psi/We) = (\kappa/We)D_{\mathrm{CCD}}\psi + (\psi/We)D_{\mathrm{CCD}}\kappa$
であるから，$\kappa$ が空間変化する一般の場合は代数的因数分解 $D_{\mathrm{CCD}}(p-\kappa\psi/We)$ は使えない．この $O(h^6)$ は各演算子の精度から直接導かれる．

\subsubsection{CSF モデル誤差との関係：実測収束次数}

ただし，残差は CSF モデルが界面を有限厚さ $\varepsilon \sim O(h)$ で近似することによる\textbf{モデル誤差 $O(\varepsilon^2) \approx O(h^2)$} が律速するため，実測の寄生流れ収束次数は $O(h^2)$ 程度となる（§\ref{sec:accuracy_summary} および第\ref{sec:verification}章の静止液滴ベンチマーク参照）．Balanced-Force 条件の「Balanced」とは「数値離散化誤差成分の相殺（$O(h^6)$ への低減）」を指し，CSF モデル誤差（$O(\varepsilon^2) \approx O(h^2)$）そのものを除去するものではない点に注意されたい．
